\chapter{Wymiar, sumy proste i przestrzenie ilorazowe}
\label{ch:lin2-wymiar-sumy-ilorazy}
\index{wymiar przestrzeni wektorowej@Wymiar przestrzeni wektorowej}
\index{suma prosta@Suma prosta}
\index{przestrzen ilorazowa@Przestrzeń ilorazowa}

\noindent\textbf{Cel rozdziału.}
W Rozdziale~\ref{ch:lin-przestrzenie-wektorowe} zdefiniowaliśmy przestrzenie
wektorowe, skonstruowaliśmy bazy i dowiedliśmy jednoznaczności współrzędnych.
Teraz zbadamy wymiar dowolnej przestrzeni, sumy podprzestrzeni, sumy proste,
przestrzenie ilorazowe oraz dopełnienia algebraiczne.
Jeżeli składnikami sumy są \emph{wektory}, zawsze dopuszczamy
wyłącznie skończenie wiele niezerowych składników, nawet gdy zbiór
indeksów jest nieskończony. Nie rozważamy granic ani zbieżności
szeregów. Natomiast wyrażenie \(\sum_{i\in I}\dim_KU_i\)
oznacza \emph{sumę kardynalną} liczb \(\dim_KU_i\);
nie jest ono sumą wektorów.

W całym rozdziale \(K\) jest ciałem przemiennym, \(0_K\neq1_K\),
a wszystkie przestrzenie wektorowe są rozpatrywane nad \(K\).
Przy twierdzeniach o dowolnych bazach nieskończonych
jawnie przyjmujemy aksjomat wyboru; twierdzenia dotyczące skończonych
wymiarów nie wymagają go.

\section{Wymiar jako liczba kardynalna}
\label{sec:lin2-wymiar-kardynalny}

Dla baz skończonych jednakową liczbę wektorów
udowodniliśmy w Rozdziale~\ref{ch:lin-przestrzenie-wektorowe}
za pomocą lematu o wymianie Steinitza.
W przypadku dowolnych baz Hamela konieczny jest dodatkowy argument:
każdy wektor korzysta tylko ze skończenie wielu współrzędnych
w danej bazie.

\begin{lemat}[Nośnik współrzędnych w bazie]
\label{lem:lin2-nosnik}
Niech \(B\subseteq V\) będzie bazą Hamela.
Każdy wektor \(v\in V\) ma dokładnie jedną rodzinę
współczynników \((a_b)_{b\in B}\) o skończonym nośniku
taką, że \(v=\sum_{b\in B}a_bb\).
Oznaczmy przez \(\operatorname{supp}_B(v)\) zbiór
tych \(b\in B\), dla których \(a_b\neq0_K\).
Dla dowolnego \(A\subseteq B\) zachodzi
\[
 \operatorname{span}(A)
 =\{v\in V:\operatorname{supp}_B(v)\subseteq A\}.
\]
\end{lemat}
\begin{proof}
Z Twierdzenia~\ref{thm:lin-wspolrzedne-hamela}
wynikają istnienie przedstawienia
oraz jednoznaczność rodziny współczynników.
Jeśli \(v\in\operatorname{span}(A)\), to jest skończoną
kombinacją wektorów z \(A\). Po dopisaniu zerowych współczynników
poza \(A\) otrzymujemy przedstawienie względem \(B\).
Jednoznaczność pokazuje, że nośnik jest podzbiorem \(A\).
W przeciwnym kierunku, jeśli nośnik współczynników
\(v\) jest zawarty w \(A\), to jego jedyny zapis
w bazie \(B\) jest skończoną kombinacją elementów \(A\).
Równość obejmuje też \(v=0_V\) oraz \(A=\varnothing\).
\end{proof}

\begin{lemat}[Suma zbiorów skończonych indeksowana zbiorem nieskończonym]
\label{lem:lin2-moce-nosnikow}
Przyjmijmy aksjomat wyboru.
Niech \(I\) będzie zbiorem nieskończonym,
a \((F_i)_{i\in I}\) rodziną zbiorów skończonych.
Wtedy
\[
 \left|\bigcup_{i\in I}F_i\right|\leq |I|.
\]
\end{lemat}
\begin{proof}
Przyjmując twierdzenie o dobrym uporządkowaniu,
ustalmy dobre uporządkowania \(I\) oraz
\(D=\bigcup_{i\in I}F_i\).
Dla \(x\in D\) niech \(i(x)\) oznacza najmniejszy indeks
taki, że \(x\in F_{i(x)}\).
Uporządkowanie \(D\) wyznacza porządek
w skończonym zbiorze \(F_{i(x)}\).
Niech \(k(x)\in\mathbb N_0\) będzie numerem \(x\)
w tym skończonym uporządkowanym zbiorze.
Odwzorowanie
\(x\mapsto(i(x),k(x))\)
jest iniekcją z \(D\) do \(I\times\mathbb N_0\):
dla ustalonego \(i\) różne elementy \(F_i\)
otrzymują różne numery.
Zatem \(|D|\leq|I\times\mathbb N_0|=|I|\),
gdzie ostatnia równość jest standardowym
twierdzeniem arytmetyki kardynalnej w ZFC
dla zbioru nieskończonego \(I\).
\end{proof}

\begin{twierdzenie}[Równoliczność dowolnych baz]
\label{thm:lin2-jednakowa-moc-baz}
Przyjmijmy aksjomat wyboru.
Jeśli \(B,C\subseteq V\) są bazami Hamela
tej samej przestrzeni wektorowej \(V\),
to \(|B|=|C|\).
\end{twierdzenie}
\begin{proof}
Jeśli \(B\) jest skończona, to
z lematu o wymianie Steinitza
(Twierdzenie~\ref{thm:lin-steinitz})
każda skończona niezależna podrodzina \(C\)
ma co najwyżej \(|B|\) elementów.
Gdyby \(C\) miała więcej elementów,
można by wybrać \(|B|+1\) różnych elementów,
otrzymując sprzeczność.
Zatem \(C\) jest skończona, a jej liczność
jest równa \(|B|\) na mocy
Wniosku~\ref{cor:lin-bazy-rowna-liczba}.
To samo rozumowanie dotyczy zamienionych baz.
Pozostaje więc przypadek, gdy obie bazy są nieskończone.

Dla \(c\in C\) niech
\(F_c=\operatorname{supp}_B(c)\), skończony zbiór.
Pokażemy, że \(B=\bigcup_{c\in C}F_c\).
Zbiór \(D=\bigcup_{c\in C}F_c\) zawiera
wszystkie wektory potrzebne do zapisania
każdego \(c\in C\).
Stąd
\(\operatorname{span}(C)\subseteq\operatorname{span}(D)\).
Ponieważ \(C\) jest bazą \(V\),
otrzymujemy \(\operatorname{span}(D)=V\).
Gdyby \(b\in B\setminus D\), to
\(\operatorname{supp}_B(b)=\{b\}\not\subseteq D\),
a więc \(b\notin\operatorname{span}(D)\)
na mocy Lematu~\ref{lem:lin2-nosnik}, sprzeczność.
Wykazaliśmy zatem \(B=D\).

Z Lematu~\ref{lem:lin2-moce-nosnikow} otrzymujemy
\(|B|\leq|C|\).
Zamieniając role baz, dostajemy \(|C|\leq|B|\).
Twierdzenie Cantora--Bernsteina daje \(|B|=|C|\).
\end{proof}

\begin{definicja}[Wymiar w dowolnym wymiarze]
\label{def:lin2-wymiar}
Zakładając aksjomat wyboru, przez
\(\dim_K V\) rozumiemy moc dowolnej bazy Hamela \(V\).
Ta liczba kardynalna nie zależy od bazy
na mocy Twierdzenia~\ref{thm:lin2-jednakowa-moc-baz}.
W szczególności \(\dim_K\{0_V\}=0\).
\end{definicja}

\begin{przyklad}[Wymiar wielomianów i ciągów]
\label{prz:lin2-przeliczalna-baza}
W przestrzeni \(K[X]\) rodzina
\((1,X,X^2,\ldots)\) jest bazą.
Także \(K^{(\mathbb N_0)}\), przestrzeń ciągów
o skończonym nośniku, ma bazę \((e_n)_{n\in\mathbb N_0}\).
Obie bazy są przeliczalnie nieskończone:
\(\dim_KK[X]=\dim_KK^{(\mathbb N_0)}=\aleph_0\).
Izomorfizm tych przestrzeni przyporządkowuje
wielomianowi ciąg jego współczynników.
Nie należy mylić przestrzeni \(K^{(\mathbb N_0)}\)
z przestrzenią wszystkich ciągów \(K^{\mathbb N_0}\):
np. \((1,1,\ldots)\) nie należy do pierwszej z nich.
\end{przyklad}

\begin{uwaga}[Rachunek wymiarów nieskończonych]
\label{rem:lin2-odejmowanie-kardynalne}
Dla skończonych wymiarów możemy odejmować liczby naturalne.
Dla liczb kardynalnych nieskończonych odejmowanie
nie jest jednoznacznie odwracalne:
na przykład \(\aleph_0+1=\aleph_0+2=\aleph_0\).
Dlatego w ogólnych wzorach wymiarowych
będziemy stosować przede wszystkim dodawanie.
\end{uwaga}

\section{Suma algebraiczna podprzestrzeni}
\label{sec:lin2-suma-algebraiczna}

\begin{definicja}[Suma podprzestrzeni]
\label{def:lin2-suma-podprzestrzeni}
Niech \((U_i)_{i\in I}\) będzie rodziną podprzestrzeni \(V\).
\emph{Sumą algebraiczną} tej rodziny nazywamy
\[
 \sum_{i\in I}U_i:=
 \left\{\sum_{i\in F}u_i:
 F\subseteq I\text{ skończony},\
 u_i\in U_i\text{ dla każdego }i\in F\right\}.
\]
Przy \(I=\varnothing\) przyjmujemy sumę pustą
równą \(\{0_V\}\).
Dla dwóch podprzestrzeni piszemy
\(U+W=\{u+w:u\in U,\ w\in W\}\).
Nie należy mylić sumy algebraicznej \(U+W\)
z \emph{sumą zbiorów} \(U\cup W\).
\end{definicja}

\begin{twierdzenie}[Najmniejsza podprzestrzeń zawierająca rodzinę]
\label{thm:lin2-suma-minimalna}
Dla każdej rodziny \((U_i)_{i\in I}\)
podprzestrzeni \(V\) zbiór
\(\sum_{i\in I}U_i\) jest podprzestrzenią
i zachodzi
\[
 \sum_{i\in I}U_i=
 \operatorname{span}\!\left(\bigcup_{i\in I}U_i\right).
\]
Jest to najmniejsza podprzestrzeń \(V\)
zawierająca wszystkie \(U_i\).
\end{twierdzenie}
\begin{proof}
Suma zawiera \(0_V\), bo dopuszczamy sumę pustą.
Niech \(x=\sum_{i\in F}u_i\),
\(y=\sum_{i\in G}w_i\), gdzie \(F,G\subseteq I\)
są skończone, \(u_i\in U_i\), \(w_i\in U_i\).
Dopiszmy wektory zerowe przy brakujących indeksach,
dzięki czemu dla \(a,b\in K\)
\[
 ax+by=\sum_{i\in F\cup G}(au_i+bw_i).
\]
Każdy składnik należy do \(U_i\),
więc suma jest podprzestrzenią.

Każda \(U_i\) zawiera się w tej sumie:
wektor \(u\in U_i\) otrzymujemy z \(F=\{i\}\).
Zatem \(\operatorname{span}(\bigcup_iU_i)\)
jest zawarta w \(\sum_iU_i\).
W przeciwną stronę każdy składnik \(u_i\)
dowolnej sumy należy do \(\bigcup_iU_i\),
więc całe wyrażenie jest jej skończoną kombinacją.
Równość i minimalność wynikają z
Twierdzenia~\ref{thm:lin-span-minimalny}.
Przy \(I=\varnothing\) obie strony są
równe \(\{0_V\}\).
\end{proof}

\begin{stwierdzenie}[Rachunek sum podprzestrzeni]
\label{prop:lin2-rachunek-sum}
Dla podprzestrzeni \(U,W,Z\subseteq V\) zachodzą
\[
 U+W=W+U,\quad
 (U+W)+Z=U+(W+Z),\quad
 U+\{0_V\}=U,\quad U+U=U.
\]
Jeżeli \(U\subseteq W\), to \(U+Z\subseteq W+Z\).
Ponadto \(U\cap W\subseteq U+W\).
\end{stwierdzenie}
\begin{proof}
Pierwszą równość daje przemienność dodawania,
drugą łączność: obie strony składają się
dokładnie z wektorów \(u+w+z\).
Kolejne dwie równości wynikają
z obecności wektora zerowego i domknięcia
podprzestrzeni na dodawanie.
Zawieranie przy \(U\subseteq W\) jest bezpośrednie.
Dla \(v\in U\cap W\) mamy
\(v=v+0_V\in U+W\).
\end{proof}

\begin{przyklad}[Suma prostych na płaszczyźnie]
\label{prz:lin2-suma-prostych}
Niech \(U=\operatorname{span}\{(1,0)\}\),
\(W=\operatorname{span}\{(1,1)\}\subseteq\mathbb R^2\).
Dla każdego \((x,y)\in\mathbb R^2\)
otrzymujemy
\((x,y)=(x-y)(1,0)+y(1,1)\),
zatem \(U+W=\mathbb R^2\).
Jeśli \(a(1,0)=b(1,1)\), to
z drugiej współrzędnej \(b=0\),
a następnie \(a=0\).
Stąd \(U\cap W=\{0\}\).
Natomiast dla
\(Z=\operatorname{span}\{(2,0)\}\)
zachodzi \(U+Z=U\), chociaż
\(\dim U=\dim Z=1\).
\end{przyklad}

\begin{twierdzenie}[Wzór wymiarowy dla dwóch podprzestrzeni]
\label{thm:lin2-wzor-wymiarowy}
Przyjmijmy aksjomat wyboru.
Jeżeli \(U,W\) są podprzestrzeniami \(V\), to
\[
 \dim_K(U+W)+\dim_K(U\cap W)=\dim_KU+\dim_KW,
\]
gdzie dodawanie jest dodawaniem liczb kardynalnych.
Jeżeli \(U\) i \(W\) są skończenie wymiarowe, to
\[
 \dim_K(U+W)=\dim_KU+\dim_KW-\dim_K(U\cap W)
\]
jako równość liczb naturalnych, bez potrzeby
stosowania aksjomatu wyboru.
\end{twierdzenie}
\begin{proof}
Wybierzmy bazę \(C\) podprzestrzeni \(U\cap W\).
Rozszerzmy ją do bazy \(C\cup A\) przestrzeni \(U\)
oraz do bazy \(C\cup D\) przestrzeni \(W\).
Możemy przyjąć \(C\cap A=C\cap D=\varnothing\).
Ponadto \(A\cap D=\varnothing\):
każdy \(x\in A\cap D\) należałby do \(U\cap W\),
więc do \(\operatorname{span}(C)\),
co przeczyłoby niezależności \(C\cup A\).

Sprawdźmy, że \(C\cup A\cup D\) jest bazą \(U+W\).
Z definicji \(U+W\), każda jego składowa
\(u\in U\) jest kombinacją \(C\cup A\),
a każda \(w\in W\) jest kombinacją \(C\cup D\).
Wobec tego \(C\cup A\cup D\) rozpina \(U+W\).

Aby sprawdzić warunek niezależności liniowej,
przypuśćmy, że skończona kombinacja
elementów \(C\cup A\cup D\) jest równa zeru.
Po zebraniu składników otrzymujemy
\(c+a+d=0_V\), gdzie
\(c\in\operatorname{span}(C)\),
\(a\in\operatorname{span}(A)\),
\(d\in\operatorname{span}(D)\).
Ponieważ \(d=-(c+a)\in U\) oraz \(d\in W\),
mamy \(d\in U\cap W=\operatorname{span}(C)\).
Zapiszmy \(d\) raz jako skończoną kombinację
elementów \(D\), a drugi raz jako skończoną kombinację
elementów \(C\). Odejmując te dwa przedstawienia,
otrzymujemy równość z warunku niezależności liniowej
zbioru \(C\cup D\). Ponieważ \(C\cap D=\varnothing\),
wszystkie współczynniki w obu przedstawieniach
muszą być zerowe, a więc \(d=0_V\).
Następnie \(c+a=0_V\), a niezależność
\(C\cup A\) wymusza zerowanie również
pozostałych współczynników.
Wykazaliśmy niezależność całej rodziny.

Wobec tego, jako równości kardynalne,
\[
 \begin{aligned}
 \dim(U+W)&=|C|+|A|+|D|,&
 \dim(U\cap W)&=|C|,\\
 \dim U&=|C|+|A|,&
 \dim W&=|C|+|D|.
 \end{aligned}
\]
Po dodaniu stron otrzymujemy pierwszy wzór.
Dla wymiarów skończonych wolno odjąć
\(|C|=\dim(U\cap W)\) od obu stron,
co daje drugi wzór. W tym przypadku
bazy i ich rozszerzenia konstruujemy
skończonym algorytmem, bez aksjomatu wyboru.
\end{proof}

\begin{przyklad}[Płaszczyzny w \(\mathbb R^3\)]
\label{prz:lin2-plaszczyzny}
Niech \(U=\{(x,y,z):z=0\}\),
\(W=\{(x,y,z):y=0\}\subseteq\mathbb R^3\).
Przestrzenie te mają bazy \((e_1,e_2)\)
oraz \((e_1,e_3)\), więc wymiar \(2\).
Przecięcie to
\(U\cap W=\operatorname{span}\{e_1\}\)
o wymiarze \(1\).
Suma zawiera \(e_1,e_2,e_3\),
które rozpinają \(\mathbb R^3\),
więc \(U+W=\mathbb R^3\).
Wzór wymiarowy przyjmuje postać \(3+1=2+2\).
\end{przyklad}


\section{Sumy proste i jednoznaczność rozkładu}
\label{sec:lin2-sumy-proste}

\begin{definicja}[Wewnętrzna suma prosta]
\label{def:lin2-wewnetrzna-suma-prosta}
Niech \((U_i)_{i\in I}\) będzie rodziną
podprzestrzeni \(V\). Jej suma algebraiczna
jest \emph{sumą prostą}, jeśli dla każdego
skończonego \(F\subseteq I\) oraz
dowolnych \(u_i\in U_i\), \(i\in F\),
z równości
\[
 \sum_{i\in F}u_i=0_V
\]
wynika \(u_i=0_V\) dla wszystkich \(i\in F\).
Taką sumę oznaczamy przez
\(\bigoplus_{i\in I}U_i\).
Zapis \(V=\bigoplus_{i\in I}U_i\)
oznacza dodatkowo, że suma jest całą przestrzenią \(V\).
Dla \(I=\varnothing\) suma jest prosta
i jest równa \(\{0_V\}\).
\end{definicja}

\begin{twierdzenie}[Kryteria sumy prostej]
\label{thm:lin2-kryteria-sumy-prostej}
Dla dowolnej rodziny \((U_i)_{i\in I}\)
podprzestrzeni \(V\) następujące warunki są równoważne:
\begin{enumerate}[label=(\roman*)]
\item suma \(\sum_{i\in I}U_i\) jest prosta;
\item każdy wektor \(v\in\sum_{i\in I}U_i\)
ma dokładnie jedno przedstawienie
\(v=\sum_{i\in I}u_i\) z \(u_i\in U_i\)
i skończonym nośnikiem \(\{i:u_i\neq0_V\}\);
\item dla każdego \(j\in I\) zachodzi
\[
 U_j\cap\left(\sum_{i\in I\setminus\{j\}}U_i\right)
 =\{0_V\}.
\]
\end{enumerate}
Gdy \(I=\varnothing\), wszystkie trzy
warunki są spełnione.
Dla dwóch podprzestrzeni kryterium redukuje
się do \(U\cap W=\{0_V\}\).
\end{twierdzenie}
\begin{proof}
(i)\(\Rightarrow\)(ii): istnienie rozkładu
wynika z definicji sumy.
Jeśli \(v=\sum_i u_i=\sum_iw_i\)
dla dwóch rodzin o skończonych nośnikach,
to \(\sum_i(u_i-w_i)=0_V\).
Suma jest skończona, a warunek prostości
wymusza \(u_i-w_i=0_V\) dla wszystkich \(i\).
Stąd oba rozkłady są identyczne.

(ii)\(\Rightarrow\)(i): zerowy wektor
posiada rozkład o wszystkich składnikach zerowych.
Jeśli \(\sum_{i\in F}u_i=0_V\),
to otrzymujemy inny rozkład z wektorami \(u_i\),
rozszerzony zerami poza \(F\).
Jednoznaczność rozkładu oznacza,
że \(u_i=0_V\) dla każdego \(i\in F\).

(i)\(\Rightarrow\)(iii): weźmy
\(v\in U_j\cap\sum_{i\neq j}U_i\).
Istnieje skończony \(F\subseteq I\setminus\{j\}\)
oraz \(u_i\in U_i\) takie, że
\(v=\sum_{i\in F}u_i\).
Równość \(v-\sum_{i\in F}u_i=0_V\)
i prostość sumy dają \(v=0_V\).
W przeciwną stronę przyjmijmy (iii)
i niech \(\sum_{i\in F}u_i=0_V\).
Dla każdego \(j\in F\) mamy
\(u_j=-\sum_{i\in F\setminus\{j\}}u_i\).
Jest to element \(U_j\)
i sumy pozostałych podprzestrzeni,
zatem z (iii) otrzymujemy \(u_j=0_V\).
Dowodzi to (i).
Przy pustym \(I\) nie ma niezerowych składników
ani dodatkowych warunków do sprawdzenia.
\end{proof}

\begin{przyklad}[Zerowe przecięcia par nie wystarczają]
\label{prz:lin2-trzy-proste}
W \(\mathbb R^2\) rozważmy proste
\[
 U_1=\operatorname{span}\{(1,0)\},\qquad
 U_2=\operatorname{span}\{(0,1)\},\qquad
 U_3=\operatorname{span}\{(1,1)\}.
\]
Każde dwie z nich mają przecięcie \(\{0\}\):
żadne z różnych podanych wektorów
nie jest skalarną wielokrotnością drugiego.
Mimo to ich suma nie jest prosta, ponieważ
\((1,0)+(0,1)-(1,1)=(0,0)\),
a wszystkie trzy składniki są niezerowe.
W kryterium dla trzech i większej liczby
podprzestrzeni nie wystarczą więc
zerowe przecięcia par.
\end{przyklad}

\begin{definicja}[Zewnętrzna suma prosta]
\label{def:lin2-zewnetrzna-suma-prosta}
Dla rodziny \((U_i)_{i\in I}\)
przestrzeni wektorowych nad tym samym ciałem
definiujemy \emph{zewnętrzną sumę prostą}
\[
 \bigoplus_{i\in I}^{\mathrm{zew}}U_i
 :=\left\{(u_i)_{i\in I}\in\prod_{i\in I}U_i:
  \{i\in I:u_i\neq0_{U_i}\}\text{ jest skończony}\right\}.
\]
Dodawanie i mnożenie przez skalary
wykonujemy współrzędna po współrzędnej.
Dla \(I=\varnothing\) otrzymujemy
jednopunktową przestrzeń zerową.
\end{definicja}

\begin{twierdzenie}[Suma zewnętrzna a suma wewnętrzna]
\label{thm:lin2-zewn-wewn}
Jeżeli wszystkie \(U_i\) są podprzestrzeniami
jednej przestrzeni \(V\), to odwzorowanie
\[
 \Sigma\colon\bigoplus_{i\in I}^{\mathrm{zew}}U_i
 \longrightarrow\sum_{i\in I}U_i,\qquad
 (u_i)_{i\in I}\longmapsto\sum_{i\in I}u_i
\]
jest liniowe i surjektywne.
Jest izomorfizmem wtedy i tylko wtedy,
gdy suma podprzestrzeni jest prosta.
\end{twierdzenie}
\begin{proof}
Rodzina o skończonym nośniku ma dobrze określoną
skończoną sumę. Produkt \(\prod_{i\in I}U_i\)
jest przestrzenią wektorową z działaniami
współrzędnymi: każdy aksjomat wynika
z odpowiedniego aksjomatu w \(U_i\),
sprawdzonego dla każdego \(i\).
Przestrzeń zewnętrzna jest jego podprzestrzenią:
nośnik kombinacji \(a(u_i)+b(w_i)\)
zawiera się w sumie mnogościowej
dwóch skończonych zbiorów nośników.
Działania są zatem zamknięte.
Dla \(a,b\in K\) zachodzi
\(\Sigma(a(u_i)+b(w_i))
=a\Sigma((u_i))+b\Sigma((w_i))\);
oba rachunki obejmują tylko skończenie wiele indeksów.
Z definicji sumy algebraicznej każdy jej element
ma przedstawienie \(\sum_i u_i\),
więc \(\Sigma\) jest surjektywne.
Jest iniektywne dokładnie wtedy,
gdy \(\Sigma((u_i))=0_V\) wymusza
\(u_i=0_V\) dla wszystkich \(i\).
Jest to właśnie warunek prostości.
\end{proof}

\begin{twierdzenie}[Baza i wymiar sumy prostej]
\label{thm:lin2-wymiar-sumy-prostej}
Przyjmijmy aksjomat wyboru.
Niech \(V=\bigoplus_{i\in I}U_i\),
a \(B_i=(b_{i,j})_{j\in J_i}\)
niech będzie bazą Hamela \(U_i\).
Rodzina wszystkich \(b_{i,j}\)
z indeksami \((i,j)\), \(i\in I\), \(j\in J_i\),
jest bazą \(V\). Stąd
\[
 \dim_KV=\sum_{i\in I}\dim_KU_i
\]
jako równość liczb kardynalnych.
Jeśli bazy \(B_i\) są już wskazane,
część twierdzenia o bazie nie wymaga
dodatkowego zastosowania aksjomatu wyboru.
\end{twierdzenie}
\begin{proof}
Każdy \(v\in V\) jest skończoną sumą
\(u_i\in U_i\). Każdy z tych wektorów
jest skończoną kombinacją elementów \(B_i\).
Łącznie \(v\) jest skończoną kombinacją
elementów zestawionej rodziny,
która zatem rozpina \(V\).

Sprawdźmy warunek niezależności liniowej.
Jeśli skończona kombinacja elementów \(B_i\)
ze współczynnikami \(a_{i,j}\) daje zero,
zgrupujmy wyrazy według pierwszego indeksu:
\[
 \sum_{i\in F}
 \left(\sum_{j\in G_i}a_{i,j}b_{i,j}\right)=0_V.
\]
Prostość sumy podprzestrzeni daje,
że każde wyrażenie w nawiasie jest zerowe.
Niezależność rodziny \(B_i\) wymusza
\(a_{i,j}=0_K\) dla wszystkich \(i,j\).
Zestawiona rodzina jest zatem bazą.
Jej indeksy tworzą rozłączną sumę
zbiorów \(\{i\}\times J_i\);
liczba kardynalna elementów takiej sumy
jest z definicji sumą kardynalną liczb \(|J_i|\).
Otrzymujemy wzór.
\end{proof}

\begin{przyklad}[Rozkład wielomianów według parzystości potęg]
\label{prz:lin2-parzyste-nieparzyste}
Niech \(E\subseteq K[X]\) składa się
z wielomianów o niezerowych współczynnikach
wyłącznie przy parzystych potęgach,
a \(O\subseteq K[X]\) analogicznie
przy nieparzystych potęgach.
Obie przestrzenie zawierają wielomian zerowy.
Każdy wielomian \(p=\sum_{j=0}^d a_jX^j\)
ma jednoznaczny rozkład \(p=p_E+p_O\),
gdzie
\[
 p_E=\sum_{\substack{0\leq j\leq d\\j\text{ parzysty}}}a_jX^j,
 \qquad
 p_O=\sum_{\substack{0\leq j\leq d\\j\text{ nieparzysty}}}a_jX^j.
\]
Jeśli wielomian należy do \(E\cap O\),
to wszystkie jego współczynniki są zerowe.
Stąd \(K[X]=E\oplus O\).
Rodziny \((1,X^2,X^4,\ldots)\) oraz
\((X,X^3,X^5,\ldots)\) są bazami odpowiednio
\(E\) i \(O\). Wszystkie trzy przestrzenie
mają wymiar \(\aleph_0\), zgodnie
z \(\aleph_0+\aleph_0=\aleph_0\).
\end{przyklad}

\section{Przestrzenie ilorazowe}
\label{sec:lin2-ilorazy}

\begin{definicja}[Warstwy i iloraz]
\label{def:lin2-iloraz}
Niech \(U\subseteq V\) będzie podprzestrzenią.
Dla \(v\in V\) przesunięciem zbioru \(U\)
o wektor \(v\) nazywamy zbiór
\[
 v+U=\{v+u:u\in U\}.
\]
Zbiór wszystkich takich warstw oznaczamy przez
\(V/U=\{v+U:v\in V\}\)
i nazywamy \emph{przestrzenią ilorazową}.
Warstwę \(v+U\) oznaczamy także przez \([v]_U\).
\end{definicja}

\begin{lemat}[Kiedy dwie warstwy są równe?]
\label{lem:lin2-warstwy}
Dla \(v,w\in V\) i podprzestrzeni \(U\subseteq V\)
następujące warunki są równoważne:
\begin{enumerate}[label=(\roman*)]
\item \(v+U=w+U\);
\item \((v+U)\cap(w+U)\neq\varnothing\);
\item \(v-w\in U\).
\end{enumerate}
Zatem dwie warstwy są równe albo rozłączne.
\end{lemat}
\begin{proof}
Jeśli warstwy są równe, ich przecięcie
zawiera \(v\), bo \(0_V\in U\).
Jeśli \(x\) należy do obu warstw, to
\(x=v+u_1=w+u_2\) dla pewnych \(u_1,u_2\in U\),
a zatem \(v-w=u_2-u_1\in U\).
Jeśli \(v-w=u_0\in U\), to
dla każdego \(u\in U\)
mamy \(v+u=w+(u_0+u)\in w+U\).
Otrzymujemy \(v+U\subseteq w+U\).
Analogicznie z \(w-v=-u_0\in U\)
dostajemy zawieranie odwrotne.
\end{proof}

\begin{twierdzenie}[Iloraz jest przestrzenią wektorową]
\label{thm:lin2-struktura-ilorazu}
Na \(V/U\) wzory
\[
 (v+U)+(w+U):=(v+w)+U,\qquad
 a(v+U):=(av)+U
\]
definiują dobrze określone działania
i nadają \(V/U\) strukturę przestrzeni
wektorowej nad \(K\).
Warstwa \(U=0_V+U\) jest zerem ilorazu.
Odwzorowanie kanoniczne
\[
 \pi_U\colon V\longrightarrow V/U,\qquad
 v\longmapsto v+U
\]
jest liniowe, surjektywne, a
\(\pi_U(v)=0_{V/U}\) dokładnie wtedy, gdy \(v\in U\).
\end{twierdzenie}
\begin{proof}
Jeśli \(v+U=v'+U\) i \(w+U=w'+U\),
to na mocy Lematu~\ref{lem:lin2-warstwy}
wektory \(v-v'\) i \(w-w'\) należą do \(U\).
Zatem
\[
 (v+w)-(v'+w')=(v-v')+(w-w')\in U,
 \qquad av-av'=a(v-v')\in U.
\]
Ponownie stosując lemat, widzimy,
że wynik dodawania warstw oraz ich
mnożenia przez skalar nie zależy
od wybranych przedstawicieli.

Aksjomaty przestrzeni wektorowej
sprawdzamy na przedstawicielach.
Na przykład
\[
 (a+b)(v+U)=((a+b)v)+U
 =(av+bv)+U=a(v+U)+b(v+U).
\]
Zera i wektora przeciwnego dostarczają
odpowiednio warstwy \(U\) i \((-v)+U\).
Łączność, przemienność, pozostałe
rozdzielności i działanie jedynki
dziedziczymy z \(V\).

Liniowość \(\pi_U\) wynika z definicji działań:
\(\pi_U(av+bw)=a\pi_U(v)+b\pi_U(w)\).
Każda warstwa ma przedstawiciela,
więc \(\pi_U\) jest surjektywne.
Ostatni warunek jest równoważnością
\(v+U=0_V+U\iff v\in U\),
udowodnioną w poprzednim lemacie.
\end{proof}

\begin{przyklad}[Warstwy poziome w \(\mathbb R^2\)]
\label{prz:lin2-warstwy-r2}
Niech \(V=\mathbb R^2\),
\(U=\mathbb R\times\{0\}\).
Wówczas dla \((a,b)\in\mathbb R^2\)
\[
 (a,b)+U=\{(x,b):x\in\mathbb R\}.
\]
Warstwa zależy tylko od liczby \(b\),
a dwa różne \(b\) dają rozłączne
proste poziome.
Przyporządkowanie \((a,b)+U\mapsto b\)
jest dobrze określoną liniową bijekcją
z \(\mathbb R^2/U\) na \(\mathbb R\).
Równość \((a,b)+U=(c,d)+U\)
zachodzi dokładnie wtedy, gdy \(b=d\).
\end{przyklad}

\begin{figure}[!htbp]
\centering
\begin{tikzpicture}[>=Latex,scale=1,every node/.style={font=\small}]
 \draw[->,gray!75] (-0.4,0)--(6,0) node[right] {\(x\)};
 \draw[->,gray!75] (0,-0.7)--(0,3.15) node[above] {\(y\)};
 \draw[very thick,manifoldblue] (-0.15,0)--(5.45,0);
 \draw[thick,manifoldred] (-0.15,1.10)--(5.45,1.10);
 \draw[thick,manifoldgreen] (-0.15,2.25)--(5.45,2.25);
 \node[manifoldblue,below] at (4.8,0) {\(U\)};
 \node[manifoldred,above] at (4.20,1.10) {\((a,b)+U\)};
 \node[manifoldgreen,above] at (4.20,2.25) {\((c,d)+U\)};
 \node[left] at (0,1.10) {\(b\)};
 \node[left] at (0,2.25) {\(d\)};
 \fill (1.0,1.10) circle (1.5pt);
 \fill (1.0,2.25) circle (1.5pt);
\end{tikzpicture}
\caption{Dwie różne warstwy względem
\(U=\mathbb R\times\{0\}\). Ich suma
mnogościowa jest podzbiorem \(\mathbb R^2\),
natomiast każda z nich jest jednym
elementem przestrzeni ilorazowej.}
\label{fig:lin2-warstwy-r2}
\end{figure}

\begin{stwierdzenie}[Przypadki skrajne ilorazu]
\label{prop:lin2-iloraz-przypadki}
Dla każdej przestrzeni wektorowej \(V\)
zachodzą naturalne izomorfizmy
\(V/\{0_V\}\cong V\) oraz \(V/V\cong\{0\}\).
\end{stwierdzenie}
\begin{proof}
Gdy \(U=\{0_V\}\), każda warstwa jest
jednoelementowym zbiorem \(v+\{0_V\}=\{v\}\),
a odwzorowanie \(\{v\}\mapsto v\)
jest liniową bijekcją.
Gdy \(U=V\), każda warstwa jest równa \(V\),
więc iloraz ma dokładnie jeden element.
Jest zatem przestrzenią zerową.
\end{proof}

\begin{twierdzenie}[Uniwersalna własność przestrzeni ilorazowej]
\label{thm:lin2-wlasnosc-ilorazu}
Niech \(V,Z\) będą przestrzeniami wektorowymi
nad \(K\), \(U\subseteq V\) podprzestrzenią,
a \(T\colon V\to Z\) odwzorowaniem liniowym,
które spełnia \(T(u)=0_Z\) dla każdego \(u\in U\).
Wówczas istnieje dokładnie jedno odwzorowanie
liniowe \(\overline T\colon V/U\to Z\) takie, że
\[
 \overline T\circ\pi_U=T.
\]
Określa je wzór \(\overline T(v+U)=T(v)\).
\end{twierdzenie}
\begin{proof}
Jeśli \(v+U=w+U\), to
\(v-w\in U\) na mocy
Lematu~\ref{lem:lin2-warstwy}.
Stąd
\(T(v)-T(w)=T(v-w)=0_Z\),
czyli \(T(v)=T(w)\).
Wzór na \(\overline T\) jest więc
dobrze określony.
Dla \(a,b\in K\) oraz \(v,w\in V\)
mamy
\[
 \overline T(a(v+U)+b(w+U))
 =\overline T((av+bw)+U)
 =T(av+bw)
 =aT(v)+bT(w).
\]
Dowodzi to liniowości.
Równość \(\overline T\circ\pi_U=T\)
zachodzi dla każdego \(v\) wprost
z definicji.
Jeśli \(S\colon V/U\to Z\) jest innym
odwzorowaniem z \(S\circ\pi_U=T\),
to dla dowolnej warstwy \(v+U\)
otrzymujemy \(S(v+U)=S(\pi_U(v))=T(v)\).
Zatem \(S=\overline T\).
\end{proof}

\section{Dopełnienia algebraiczne i wymiar ilorazu}
\label{sec:lin2-dopelnienia}

\begin{definicja}[Dopełnienie podprzestrzeni]
\label{def:lin2-dopelnienie}
Podprzestrzeń \(W\subseteq V\) nazywamy
\emph{dopełnieniem algebraicznym}
podprzestrzeni \(U\subseteq V\), jeśli
\(V=U\oplus W\); czyli \(U+W=V\)
oraz \(U\cap W=\{0_V\}\).
\end{definicja}

\begin{twierdzenie}[Istnienie dopełnienia algebraicznego]
\label{thm:lin2-istnienie-dopelnienia}
Przyjmijmy aksjomat wyboru.
Niech \(U\subseteq V\) będzie podprzestrzenią.
Wówczas istnieje podprzestrzeń \(W\subseteq V\)
taka, że \(V=U\oplus W\).
Jeżeli \(V\) jest skończenie wymiarowa,
twierdzenie nie wymaga aksjomatu wyboru.
\end{twierdzenie}
\begin{proof}
Weźmy bazę Hamela \(B_U\) przestrzeni \(U\).
Z twierdzenia o rozszerzaniu bazy
(Twierdzenie~\ref{thm:lin-istnienie-bazy-zorn})
istnieje baza \(B\) całej \(V\)
taka, że \(B_U\subseteq B\).
Niech \(C=B\setminus B_U\) oraz
\(W=\operatorname{span}(C)\).
Ponieważ \(B=B_U\sqcup C\) rozpina \(V\),
mamy \(U+W=V\).
Jeśli \(v\in U\cap W\), to
\[
 v=\sum_{b\in F}a_bb=\sum_{c\in G}d_cc,
 \quad F\subseteq B_U,\ G\subseteq C\text{ skończone}.
\]
Po przeniesieniu składników na jedną stronę
otrzymujemy równość, do której stosujemy
warunek niezależności liniowej bazy \(B\).
Wszystkie współczynniki są więc zerowe,
czyli \(v=0_V\).
Zatem \(U\cap W=\{0_V\}\), co dowodzi
\(V=U\oplus W\).
Jeśli \(V\) ma skończoną bazę, także \(U\)
ma skończoną bazę, a jej rozszerzenie
można wykonać skończoną procedurą
z Twierdzenia~\ref{thm:lin-uzupelnianie-skonczone}.
\end{proof}

\begin{twierdzenie}[Iloraz i dopełnienie]
\label{thm:lin2-iloraz-dopelnienie}
Jeżeli \(V=U\oplus W\), to
\[
 \Theta\colon W\longrightarrow V/U,\qquad
 w\longmapsto w+U
\]
jest izomorfizmem przestrzeni wektorowych.
Przyjmując aksjomat wyboru,
dla każdej podprzestrzeni \(U\subseteq V\)
otrzymujemy równość kardynalną
\[
 \dim_KV=\dim_KU+\dim_K(V/U).
\]
W przypadku skończenie wymiarowym
wynika z niej
\(\dim_K(V/U)=\dim_KV-\dim_KU\).
\end{twierdzenie}
\begin{proof}
Liniowość \(\Theta\) jest ograniczeniem
odwzorowania ilorazowego \(\pi_U\).
Jeżeli \(\Theta(w)=0_{V/U}\),
to \(w\in U\) na mocy
Twierdzenia~\ref{thm:lin2-struktura-ilorazu}.
Ponieważ \(w\in W\) i \(U\cap W=\{0_V\}\),
otrzymujemy \(w=0_V\), więc \(\Theta\)
jest iniektywne.
Niech \(v+U\in V/U\). Z rozkładu
\(v=u+w\), gdzie \(u\in U\), \(w\in W\),
otrzymujemy \(v+U=w+U=\Theta(w)\),
co dowodzi surjektywności.

Z Twierdzenia~\ref{thm:lin2-istnienie-dopelnienia}
wybieramy algebraiczne dopełnienie \(W\).
Izomorfizm \(\Theta\) przenosi każdą bazę
\((w_j)_{j\in J}\) przestrzeni \(W\)
na bazę \((\Theta(w_j))_{j\in J}\) ilorazu:
jego surjektywność daje rozpinanie,
gdyż każda warstwa jest obrazem pewnego
wektora z \(W\), a ten jest skończoną
kombinacją elementów bazy.
Jeśli
\(\sum_{j\in F}a_j\Theta(w_j)=0_{V/U}\),
to z liniowości otrzymujemy
\(\Theta(\sum_{j\in F}a_jw_j)=0_{V/U}\).
Iniektywność \(\Theta\) daje
\(\sum_{j\in F}a_jw_j=0_V\),
a niezależność rodziny \((w_j)_{j\in J}\)
wymusza \(a_j=0_K\) dla każdego \(j\in F\).
Tak więc rodzina obrazów jest bazą
indeksowaną przez ten sam zbiór \(J\).
Stąd \(\dim_KW=\dim_K(V/U)\).
Stosując następnie
Twierdzenie~\ref{thm:lin2-wymiar-sumy-prostej},
otrzymujemy
\[
 \dim_KV=\dim_KU+\dim_KW
 =\dim_KU+\dim_K(V/U).
\]
Jeżeli \(V\) jest skończenie wymiarowa,
bazy \(U\) i \(W\) otrzymujemy z twierdzeń
skończenie wymiarowych bez aksjomatu wyboru.
Połączenie tych baz jest bazą \(V=U\oplus W\):
rozpinanie wynika z równości \(U+W=V\),
a warunek niezależności liniowej
z \(U\cap W=\{0_V\}\) i niezależności
obu baz. Zatem liczba jej wektorów jest
sumą liczb wektorów obu baz.
Możemy więc w tym przypadku odjąć
\(\dim_KU\) i otrzymać wskazany wzór,
nie odwołując się do aksjomatu wyboru.
\end{proof}

\begin{przyklad}[Niekanoniczność dopełnienia]
\label{prz:lin2-dopelnienia}
W \(\mathbb R^2\) niech \(U=\mathbb R\times\{0\}\).
Zarówno \(W_1=\operatorname{span}\{(0,1)\}\),
jak i \(W_2=\operatorname{span}\{(1,1)\}\)
są dopełnieniami \(U\).
Istotnie, z równości
\(a(1,0)+b(0,1)=0\) wynika \(a=b=0\),
podobnie z
\(a(1,0)+b(1,1)=(a+b,b)=0\)
wynika \(a=b=0\).
W obu przypadkach generatory rozpinają
\(\mathbb R^2\), bo
\((x,y)=x(1,0)+y(0,1)
=(x-y)(1,0)+y(1,1)\).
Podprzestrzenie \(W_1\) i \(W_2\)
są różne, choć obie są izomorficzne
z \(\mathbb R^2/U\). Oznacza to, że
\emph{dopełnienie algebraiczne nie jest
jednoznacznie wyznaczone} przez samo zawieranie
\(U\subseteq V\). Nie utożsamiamy tego
spostrzeżenia z osobnym twierdzeniem o braku
konstrukcji naturalnej: samo wskazanie dwóch
różnych dopełnień nie wystarczałoby
do udowodnienia takiej tezy.
\end{przyklad}

\begin{twierdzenie}[Twierdzenie o ilorazie sumy]
\label{thm:lin2-iloraz-sumy}
Niech \(U,W\subseteq V\) będą podprzestrzeniami.
Wtedy istnieje naturalny izomorfizm
\[
 U/(U\cap W)\ \cong\ (U+W)/W,
 \qquad u+(U\cap W)\longmapsto u+W.
\]
\end{twierdzenie}
\begin{proof}
Określmy odwzorowanie
\(T\colon U\to (U+W)/W\),
\(T(u)=u+W\).
Jest ono liniowe, jako ograniczenie
kanonicznego odwzorowania ilorazowego
\(U+W\to(U+W)/W\).
Ponadto \(T(u)=0\) dokładnie wtedy,
gdy \(u\in W\); ponieważ \(u\in U\),
jest to równoważne \(u\in U\cap W\).
W szczególności \(T\) jest zerowe
na \(U\cap W\).
Z własności uniwersalnej
(Twierdzenie~\ref{thm:lin2-wlasnosc-ilorazu})
otrzymujemy liniowe odwzorowanie
\(\overline T\colon U/(U\cap W)\to (U+W)/W\)
dane wzorem z tezy.

Jeśli \(\overline T(u+(U\cap W))=0\),
to \(u\in W\), a więc \(u\in U\cap W\).
Oznacza to, że wyjściowa warstwa jest zerowa,
czyli \(\overline T\) jest iniektywne.
Jeśli \(x+W\in(U+W)/W\), to
\(x=u+w\) dla pewnych \(u\in U,w\in W\).
Stąd \(x+W=u+W=\overline T(u+(U\cap W))\).
Odwzorowanie jest więc surjektywne
i stanowi izomorfizm.
\end{proof}

\begin{przyklad}[Iloraz wielomianów przez wielomiany stałe]
\label{prz:lin2-iloraz-wielomiany}
Niech \(V=K[X]\), \(U=K[X]_{\leq0}\).
Podprzestrzeń \(U\) składa się ze stałych wielomianów,
a \(W=\{p\in K[X]:p(0)=0_K\}\)
ze wszystkich wielomianów o zerowym wyrazie stałym.
Dla każdego \(p(X)=a_0+\sum_{j=1}^d a_jX^j\)
(z umową, że suma jest pusta przy \(d=0\))
mamy \(p=a_0+q\), gdzie
\(q(X)=\sum_{j=1}^d a_jX^j\in W\). Wspólnym elementem
\(U\cap W\) jest tylko wielomian zerowy,
zatem \(K[X]=U\oplus W\).
Iloraz \(K[X]/K[X]_{\leq0}\) jest więc
izomorficzny z \(W\), która ma bazę
\((X,X^2,X^3,\ldots)\).
Iloraz nadal ma wymiar \(\aleph_0\),
tak jak przestrzeń wyjściowa.
\end{przyklad}

\section{Podsumowanie}
\label{sec:lin2-podsumowanie}

Suma zbiorów \(U\cup W\) i suma algebraiczna
podprzestrzeni \(U+W\) są różnymi konstrukcjami.
Suma prosta oznacza jednoznaczny rozkład na składniki
należące do odpowiednich podprzestrzeni.
W przestrzeni ilorazowej \(V/U\) utożsamiamy wektory,
których różnica należy do \(U\); otrzymane warstwy
są wektorami nowej przestrzeni.
Przy założeniu aksjomatu wyboru każda podprzestrzeń
posiada dopełnienie algebraiczne, a wymiar
zachowuje się zgodnie z dodawaniem kardynalnym:
\[
 \dim_K(U+W)+\dim_K(U\cap W)=\dim_KU+\dim_KW,
 \qquad
 \dim_KV=\dim_KU+\dim_K(V/U).
\]
Przy wymiarach skończonych wolno z tych równości
odjąć odpowiednie składniki, natomiast
w przypadku nieskończonym nie wolno tego robić
bez dodatkowych warunków.
